\clearpage\begingroup
\hypertarget{native-310-lower}{}\HistoricalHead{305}{R}{BERICHTE UND ANTRÄGE ZUR ANSCHAFFUNG VON INSTRUMENTEN.}\HistoricalFoot{XI1.}{39}
\begin{center}Unterthänigster Bericht\par\smallskip
Ueber meine Reise nach München und Benedictbeurn\par
wegen der astronomischen Instrumente.\end{center}

\noindent\hspace*{\parindent}In Folge des mir unter dem 28. März d. J. [1816] ertheilten höhern Auftrages habe ich die Reise nach München und Benedictbeurn am 18. April angetreten, mich mit den dortigen Instituten zur Verfertigung astronomischer Instrumente genau bekannt gemacht, und besonders mit dem Salinenrath \textsc{von Reichenbach} wegen derjenigen Instrumente, deren Bedürfniss für die neue Sternwarte am dringendsten sein wird, eine zwar nur provisorische, aber umständliche Rücksprache genommen. Ich habe jetzt die Ehre, hier dem hohen Curatorium über das Wesentliche unterthänigen Bericht abzustatten.

Vor etwa 10 oder 12 Jahren bildete sich in München das Institut zur Verfertigung astronomischer Instrumente, welches mit Recht eine so grosse Celebrität erhalten, und beinahe allen Sternwarten des festen Landes Instrumente geliefert hat, die den besten englischen von den Kennern gleichgesetzt und in mehr als Einer Rücksicht selbst vorgezogen werden. Die Seele dieses Institutes war der damalige k. bayrische Artilleriehauptmann, jetzt Salinenrath \textsc{von Reichenbach}, der mit seinen grossen Talenten für die Mechanik feine theoretische Einsichten verbindet, und sich auch als Schriftsteller (über die eisernen Brücken) bekannt gemacht hat. Die neue von ihm erfundene höchst sinnreiche Theilungsmethode würde allein schon allen seinen Instrumenten einen entschiedenen Vorzug geben, wenn sie sich auch nicht durch eine Menge anderer eigenthümlicher Vollkommenheiten empföhlen. In Verbindung mit ihm trat der Geheime Rath \textsc{von Utzschneider}, welcher einen Theil seines
\clearpage\SourceLeaf{311}\HistoricalHead{306}{L}{PRAKTISCHE UND SPHÄRISCHE ASTRONOMIE.}
\noindent grossen Vermögens zur Unterstützung des Instituts verwandte, und besonders in der ihm gehörenden vormaligen Abtei Benedictbeurn (16 Stunden jenseits von München) für den optischen Theil der Instrumente sehr ins Grosse gehende Anlagen machte. In der dortigen Glasfabrik wird, nach langen Versuchen, wobei keine Kosten gespart wurden, sowohl das Crownglas als das Flintglas jetzt in einer bisher nie erreichten Vollkommenheit verfertigt; die optischen Gläser werden nach eigenthümlichen Methoden bearbeitet, und die dort verfertigten achromatischen Fernröhre lassen die englischen jetzt hinter sich zurück, wie ich weiterhin noch näher anzeigen werde. Endlich trat dem Verein noch ein geschickter Mechaniker in München Namens \textsc{Liebherr} bei, welcher besonders die Verfertigung der astronomischen Uhren besorgte. Alle diese Anstalten bildeten Ein Ganzes, und die sämmtlichen daraus hervorgegangenen Instrumente sind mit der Firma: \textsc{Reichenbach, Utzschneider} und \textsc{Liebherr} bezeichnet.

Seit einigen Jahren ist nun zwar diese Verbindung wieder aufgehoben: allein die drei Werkstätten, obgleich nicht mehr unter Einer Verwaltung stehend, liefern fortwährend dieselben Arbeiten, und unterstützen dabei einander wechselseitig. Der Salinenrath \textsc{von Reichenbach} hat auf die grössern und feinern astronomischen Messwerkzeuge nach allen Theilen Europens so viele Bestellungen, dass seine grosse mit zahlreichen Arbeiten besetzte Werkstatt sie kaum bestreiten kann, zumal, da er die delicatesten Arbeiten dabei eigenhändig ausführt, wozu er, da der Maschinenbau im Grossen bei den bayrischen Salinen ihn immer sehr beschäftigt, nur einen Theil seiner Zeit verwenden kann. — Das optische Institut in Benedictbeurn steht jetzt für sich unter der unmittelbaren Leitung eines sehr talentvollen und thätigen Mannes, \textsc{Fraunhofer}, und liefert die Fernröhre, so wie alle astronomischen Instrumente, bei welchen die Fernröhre die Hauptsache sind. — \textsc{Liebherr} macht sein Hauptgeschäft aus der Verfertigung astronomischer Uhren, an denen er neuerdings eine sinnreiche ihm eigenthümliche Verbesserung angebracht hat.

Nach dieser allgemeinen Uebersicht komme ich auf die Bedürfnisse unsrer neuen Sternwarte. Ich darf mich dabei auf meine frühern unterthänigen Darstellungen beziehen, wo ich ausführlich entwickelt habe, dass, ganz dem ersten ursprünglichen Plane von 1803 gemäss, nachdem nur das \textsc{Repsold}sche Instrument (über dessen baldige Ablieferung ich posttäglich bestimmte Nachricht
\clearpage\SourceLeaf{312}\HistoricalHead{307}{R}{BERICHTE UND ANTRÄGE ZUR ANSCHAFFUNG VON INSTRUMENTEN.}\HistoricalFoot{}{39\textsuperscript{*}}
\noindent erwarte) an die Stelle des kleinern Passageinstruments und der Renovation des Mauerquadranten getreten ist, das wesentliche Bedürfniss noch in folgendem besteht:
\begin{center}\begingroup\fontsize{9.5}{12}\selectfont
\(\text{Fixe Instrumente}\left\{\begin{array}{@{}l@{\;}l@{}}
\begin{tabular}{@{}l@{}}im westlichen\\Observationszimmer\end{tabular}&\left\{\begin{tabular}{@{}rl@{}}I.&Der Meridiankreis.\\II.&Das Passageinstrument.\end{tabular}\right.\\[3pt]
\text{in der Rotunde}&\begin{tabular}{@{}rp{.39\textwidth}@{}}III.&\textls[120]{Entweder} ein Aequatoreal \textls[120]{oder} ein grosses parallatisch aufzustellendes achromatisches Fernrohr\end{tabular}
\end{array}\right.\)\par\smallskip
und dazu . . . . . . . . . . . . . . .\quad IV. Zwei astronomische Pendeluhren.
\par\endgroup\end{center}

Ich werde mich nun über diese vier Artikel umständlich erklären.
\begin{center}\textls[120]{I. Der Meridiankreis.}\end{center}
\noindent\hspace*{\parindent}Bei dem gegenwärtigen Zustande der praktischen Astronomie hat man dieses Instrument gewissermaassen als das wichtigste auf einer Sternwarte des ersten Ranges zu betrachten, und da man dergleichen Kreise seit 30 Jahren von der mannichfaltigsten Art gebauet hat, so war es ein Hauptaugenmerk meiner Reise, über die vortheilhafteste Einrichtung mit \textsc{Reichenbach} Rücksprache zu nehmen. Ich wusste bereits vorher, aus Briefen von \textsc{Reichenbach}, dass dieser jetzt damit umgeht, den astronomischen Kreisen eine ganz neue Einrichtung zu geben. Ich fand in München bereits vier solcher neuer Kreise in Arbeit, von denen drei für die Sternwarten in Ofen, Turin und Warschau bestellt sind, der vierte ohne jedoch noch bestimmt bestellt zu sein, für eine andere Sternwarte gewünscht wird. Da ich an diesen Kreisen während meines Aufenthalts in München täglich habe arbeiten sehen, die Zeichnungen davon immer vor Augen gehabt und mich mit dem Künstler über ihre Eigenthümlichkeiten besprochen habe, so kann ich aus Ueberzeugung sagen, dass diese Gattung von Kreisen, von einem Künstler wie \textsc{Reichenbach} ausgeführt, meiner Meinung nach keiner andern nachsteht, vielmehr wesentliche eigenthümliche Vorzüge hat. Der Preis eines solchen Meridiankreises ist 3500 Gulden bayrisch (24 fl.-Fuss). Das Fundament für den Meridiankreis auf der neuen Sternwarte, welches schon seit 1803 liegt, deutet zwar darauf hin, dass damals die Absicht gewesen sein mag, einen solchen Kreis zu haben wie den \textsc{Ramsden}schen, mit welchem \textsc{Piazzi} in Palermo beobachtet, und welcher 450 Guinees gekostet hat, und gegenwärtig in England wahrscheinlich noch beträchtlich
\clearpage\SourceLeaf{313}\HistoricalHead{308}{L}{PRAKTISCHE UND SPHÄRISCHE ASTRONOMIE.}
\noindent höher zu stehen kommen würde. Inzwischen kann jenes Fundament, wie ich mich durch genaue Ausmessung jetzt versichert habe, auch eben so gut für einen Kreis von der neuen \textsc{Reichenbach}schen Einrichtung dienen; und ich halte mich überzeugt, dass die neue Sternwarte, wenn sie einen solchen Kreis besitzen wird, nicht bloss mit der Palermer Sternwarte, sondern auch mit der Greenwicher \textls[120]{rühmlich} wird wetteifern können, obgleich der \textsc{Troughton}sche Meridiankreis, welchen letztere besitzt, einer Privatnachricht zu Folge 1800 Guinees gekostet hat.

Ich trage daher darauf an, einen Meridiankreis von der neuen Einrichtung bei H. \textsc{von Reichenbach} sofort bestellen zu dürfen, wo hoffentlich der vierte noch nicht bestimmt bestellte und schon ziemlich vorgerückte noch zu erhalten stehen wird. Aber auch in diesem günstigen Falle wird bis zur Vollendung noch wenigstens Ein Jahr gewartet werden müssen, und noch viel länger, wenn über den vierten schon anderweitig disponirt und ein fünfter erst von neuem anzufangen wäre.
\begin{center}\textls[120]{II. Das Passageinstrument.}\end{center}
\noindent\hspace*{\parindent}Zwei Passageinstrumente von \textsc{Reichenbach} habe ich auf meiner Reise zu sehen Gelegenheit gehabt, ein 6füssiges in München, welches der dortigen Sternwarte gehört, und ein kleineres auf der Sternwarte des Pater \textsc{Placidus Heinrich} in Regensburg; an beiden ist die Arbeit des Meisters würdig. Der Preis eines sechsfüssigen Passageinstrumentes, dergleichen \textsc{Reichenbach} auch nach Ofen, Mayland, Neapel und Pisa geliefert hat, ist 2100 fl., und ich glaube behaupten zu können, dass mit einem solchen Instrumente eben so viel geleistet werden kann, wie mit dem 8füssigen von \textsc{Ramsden} auf der Seeberger Sternwarte, welches ich genau kenne, und welches viel mehr gekostet hat. Uebrigens habe ich mit \textsc{Reichenbach} noch verschiedene nicht unwichtige an dem Passageinstrumente anzubringende Verbesserungen verabredet. \textsc{Reichenbach} erklärte sich auch bereit, Passageinstrumente von grösserer Dimension als sechsfüssig, auf Verlangen zu liefern; allein aus andern wichtigen Rücksichten hält er selbst es nicht für rathsam, über diese Grösse hinauszugehen. Da ein solches Passageinstrument gleichfalls nicht unter Einem Jahre nach der Bestellung geliefert werden kann, so trage ich unterthänigst darauf an, diese Bestellung jetzt machen zu dürfen.
\clearpage\SourceLeaf{314}\HistoricalHead{309}{R}{BERICHTE UND ANTRÄGE ZUR ANSCHAFFUNG VON INSTRUMENTEN.}
\begin{center}\textls[120]{III. Aequatoreal — oder parallatisch montirtes achromatisches}\par
\textls[120]{Fernrohr.}\end{center}
\noindent\hspace*{\parindent}Der Bau der neuen Sternwarte ist so, dass nur das Eine \textls[120]{oder} das andere von diesen beiden Instrumenten wird aufgestellt werden können, und zwar in der Rotunde mit dem Drehdach: bei der wahrhaft unvergleichlichen Vortrefflichkeit, welche beide Arten von Instrumenten, nemlich die Aequatoreale von \textsc{Reichenbach} und die grossen Achromate von \textsc{Fraunhofer} haben, ist es in der That schwer, sich zu einer Wahl zu entschliessen. Glücklicherweise haben beide Instrumente gewissermaassen ähnliche Bestimmungen, nur dass bei dem Aequatoreal mehr das Messen, bei dem parallatisch montirten Achromat mehr das Sehen die Hauptsache ausmacht. Unsrer Sternwarte würde eine Hauptzierde fehlen, wenn sie nicht entweder das eine oder das andere hätte; allein welches von beiden auch gewählt werde, so wird die Entbehrung des andern nicht zu sehr empfunden werden. Beide Instrumente gehören übrigens zu den kostbarsten: wenn dagegen in Rücksicht der Dringlichkeit eine Rangordnung gemacht und die Anschaffung der Instrumente nicht auf Einmal geschehen soll, so leidet es keinen Zweifel, dass dieses Nr. III dem Meridiankreis und dem Passageinstrument nachstehen müsse. Das Aequatoreal von \textsc{Reichenbach}, welches der Münchner Sternwarte gehört, habe ich dort genau kennen gelernt: es ist ein Meisterstück der Kunst; ähnliche Instrumente hat H. \textsc{von Reichenbach} nach Ofen, Neapel und Pisa geliefert. Der Preis ist, wenn ich nicht irre, 3500 Gulden. — Einen grossen Achromat von 10 Fuss Brennweite, \(7\frac14\) Zoll Oeffnung, parallatisch montirt und vermittelst eines Uhrwerks der täglichen Bewegung der Himmelskörper von selbst folgend, hat das optische Institut in Benedictbeurn nach Neapel geliefert, und dasselbe hat 4500 Gulden gekostet: noch niemals war bisher ein Refraktor von dieser Grösse verfertigt worden. Allein seitdem ist man noch weiter fortgeschritten: ich fand in Benedictbeurn die Gläser zu einem riesenhaften Achromat von 15 Fuss Brennweite und 9 Zoll Oeffnung vollendet, es fehlte aber noch an der Aufstellung, so dass sich die Wirkung am Himmel noch nicht beobachten liess. Der Geheime Rath \textsc{von Utzschneider} äusserte seine Bereitwilligkeit, dieses ganz einzige Instrument, wenn es vollendet sein würde, auf einige Zeit auf die hiesige Sternwarte zu schicken, ohne dass diese eine Verbindlichkeit übernähme,
\clearpage\SourceLeaf{315}\HistoricalHead{310}{L}{PRAKTISCHE UND SPHÄRISCHE ASTRONOMIE.}
\noindent sondern nur, damit erprobt würde, was sich durch die optische Kunst leisten lasse, deren möglichste Vervollkommnung, ohne Rücksicht auf Gewinn, die Tendenz des Instituts in Benedictbeurn ist.

Bei so bewandten Umständen möchte ich vorschlagen, die Bestellung des in der Rotunde aufzustellenden Instruments, ja selbst noch die Entscheidung zwischen Aequatoreal oder parallatischem Refraktor, noch eine Zeitlang zu verschieben, bis erst astronomische Erfahrungen den Werth der grossen Verbesserungen der optischen Werkzeuge in Benedictbeurn noch mehr zu würdigen erlauben.
\begin{center}\textls[120]{IV. Astronomische Uhren.}\end{center}
\noindent\hspace*{\parindent}Die neue Sternwarte bedarf wegen der Ausdehnung des Gebäudes dreier Pendeluhren, wovon eine im westlichen, eine im östlichen Observationszimmer und eine in der Rotunde nöthig sein wird. Es brauchen jedoch nur zwei neue angeschafft zu werden, indem die \textsc{Shelton}sche, obgleich schon beinahe ein halbes Jahrhundert auf der alten Sternwarte befindlich, noch immer eine sehr gute Uhr und für den zweiten Platz auf der neuen Sternwarte vollkommen zureichend ist. Von den beiden neuen Uhren müsste wenigstens die Eine von höchster Vollkommenheit sein, während an die andere nicht ganz eben so hohe Forderungen gemacht zu werden brauchen. Der Mechaniker \textsc{Liebherr}, welcher auch theils nach Neapel, theils nach Ofen, theils behuf der astronomischen Beobachtungen bei der grossen Vermessung von Bayern astronomische Uhren geliefert hat\begingroup\renewcommand{\thefootnote}{*)}\footnote{Letztere ist mir von dem Astronomen H. \textsc{Soldner}, der stets mit dieser Uhr beobachtet hat, sehr gerühmt.}\endgroup, macht sich anheischig, eine astronomische Pendeluhr in grösster Vollkommenheit etwa 6 Monat nach der Bestellung zu liefern. Das einzige, was dabei Bedenken machen könnte, ist der ausserordentlich geringe Preis, welchen er fordert, nemlich 25 Carolin oder 275 Gulden, kaum ein Viertel von dem, was eine solche Uhr in England kosten würde — wenn man nicht schon gewohnt wäre, alle Preise dort viel niedriger zu finden als auswärts. Da übrigens die dritte Uhr nicht eher nöthig sein wird, als bis in der Rotunde ein Instrument stehen wird, so trage ich unterthänigst darauf an, vorerst nur Eine Astronomische Uhr bei \textsc{Liebherr} bestellen zu dürfen. Gesetzt auch, dass dieselbe noch etwas zu wünschen übrig
\clearpage\SourceLeaf{316}\HistoricalHead{311}{R}{BERICHTE UND ANTRÄGE ZUR ANSCHAFFUNG VON INSTRUMENTEN.}
\noindent liesse (welches zu erwarten ich übrigens gar keinen Grund habe), so würde doch nichts verloren sein; denn selbst die dritte Uhr wüsste ich anderswoher für einen eben so billigen Preis nicht zu erhalten\begingroup\renewcommand{\thefootnote}{*)}\footnote{In dem zu Anfang des Jahres 1815 eingereichten P. M. über die für die neue Sternwarte nothigen Instrumente [siehe oben S. 295] hatte ich als möglich angegeben, dass vielleicht nur Eine neue Uhr anzuschaffen wäre, wenn nämlich unter den aus Lilienthal zu erwartenden Uhren eine wäre, welche sich für den dritten Platz qualificirte. Allein ich habe seitdem erfahren, dass dies nicht der Fall ist, indem von den erst noch abzuliefernden Uhren aus Lilienthal keine mit Compensation versehen ist, mithin dieselben nur zu einem untergeordneten Gebrauch werden angewandt werden können.}\endgroup.

Fasse ich also die Resultate dieses Berichtes zusammen, so geht mein unterthänigster Antrag und Bitte dahin:\par
\hspace*{\parindent}dass ich autorisirt werden möge, sofort zu bestellen
\begin{center}\begin{tabular}{@{}rp{.85\textwidth}@{}}
I.&den Meridiankreis bei dem Salinenrath \textsc{von Reichenbach}.\\
II.&Ein sechsfüssiges Passageinstrument bei demselben.\\
IV.&Eine der beiden noch nöthigen Uhren bei \textsc{Liebherr}.
\end{tabular}\end{center}

Dass dagegen noch etwas verschoben werde die Auswahl und Bestellung von Nro. III, des Instruments für die Rotunde, so wie auch die Bestellung der zweiten noch nöthigen Uhr.

Bei bedeutenden Bestellungen von Instrumenten pflegen sonst die Künstler die Vorausbezahlung eines Theils des Preises zu verlangen. Obgleich eine solche Forderung an sich nicht unbillig ist, da wenige Künstler im Stande sind, so grosse Vorschüsse aus eignen Mitteln zu machen, so scheint doch der Salinenrath \textsc{von Reichenbach} diese Forderung nicht zu machen, sondern die Zahlung erst bei Ablieferung der Instrumente zu erwarten.
\begin{center}\rule{.18\textwidth}{.35pt}\end{center}
\noindent\hspace*{\parindent}Die Specification der unmittelbaren Reisekosten habe ich die Ehre hier beizufügen. Die Dauer der Reise hat zwar 36 Tage betragen, wovon 8 Tage auf die Hinreise nach München über Mühlhausen, Gotha, Meiningen, Würzburg, Anspach, Nördlingen und Augsburg, 12 Tage auf den Aufenthalt in München, die Reise nach Benedictbeurn, den Aufenthalt daselbst und die Rückreise nach München, und 16 Tage auf die Rückreise von München über Reichenhall, Landshut, Regensburg und Nürnberg kamen. Allein da dieser Umweg auf der Rückreise und einige Tage Aufenthalt auf der Seeberger Sternwarte zu dem unmittelbaren Zweck der Reise nicht nothwendig waren,
\clearpage\SourceLeaf{317}\HistoricalHead{312}{L}{PRAKTISCHE UND SPHÄRISCHE ASTRONOMIE.}
\noindent so habe ich die Rückreise auf demselben Wege wie die Hinreise berechnet, wo dann die ganze Dauer nur 28 Tage betragen haben würde.
\begin{flushright}Unterthänigst\end{flushright}
Göttingen, den 5. Junius 1816.
\begin{flushright}C. F. \textsc{Gauss}.\end{flushright}
\bigskip
\begin{center}\textbf{BEMERKUNGEN.}\end{center}
\begingroup\small
\noindent\hspace*{\parindent}Da der Bau der neuen Göttinger Sternwarte sich infolge der politischen Verhältnisse über viele Jahre hingezogen hatte, verzögerte sich auch die Herstellung der neuen Instrumente, die dort Aufstellung finden sollten. Auch mussten die zur Anschaffung erforderlichen Mittel auf mehrere Jahre vertheilt werden, so dass \textsc{Gauss} erst verhältnissmässig spät in den Besitz der Instrumente gelangte. Die Originale der hier abgedruckten Berichte befinden sich im Kuratorium der Göttinger Universität.
\begin{flushright}\textsc{Brendel}.\end{flushright}\par\endgroup
\begin{center}\rule{.24\textwidth}{.35pt}\end{center}
\clearpage\endgroup
